% GENERATED FILE. Edit canonical lessons, then run npm run generate:books.
% canonical-source-hash: fnv1a64:3e7e71bf932a95f5
% canonical-lessons: KA-C185-notu, KA-C185-nanya, KA-C185-riyayiti, KA-C185-kilo, KA-C185-pars

\chapter{Banknote, Coin, Discount, Kilogram, Purse}
\label{ch:ka-banknote-coin-discount-kilogram-purse}
\hlchaptermodality{185}

\begin{chapteropening}
\textbf{By the end of this chapter:} \emph{I can say a banknote, a coin, a discount, a kilogram and a purse.}

Complete the last lesson of chapter 185: I can say a banknote, a coin, a discount, a kilogram and a purse.
\end{chapteropening}

\section[\texorpdfstring{\kn{ನೋಟು} (nōṭu)}{ನೋಟು (nōṭu)}]{\kn{ನೋಟು} (nōṭu) — a banknote}

\label{lesson:KA-C185-notu}



\begin{quote}
Before the new one: say the Kannada for leafy greens, then the Kannada for tender coconut water.
\end{quote}

\subsection*{You'll want to know: \kn{ನೋಟು}}
\textbf{\kn{ನೋಟು}} — \emph{nōṭu} — \textquotedblleft{}a banknote\textquotedblright{}.

\begin{grammarlens}[title={What you've built}]
One more word for this chapter.
\end{grammarlens}

\subsection*{Your turn}
\begin{itemize}
\raggedright
  \item \emph{Say it:} \emph{nōṭu}
  \item \emph{Say it:} \emph{nōṭu}, once more
  \item \emph{Say it:} say \emph{nōṭu}
  \item \emph{Recall:} say the Kannada for leafy greens, then the Kannada for tender coconut water, then say \emph{nōṭu} again
\end{itemize}

\begin{tcolorbox}[breakable,colback=teal!4,colframe=teal!35!black,title={Before you move on}]
What does \kn{ನೋಟು} mean? (\textquotedblleft{}A banknote\textquotedblright{}.) Say it once more.
\end{tcolorbox}
\section[\texorpdfstring{\kn{ನಾಣ್ಯ} (nāṇya)}{ನಾಣ್ಯ (nāṇya)}]{\kn{ನಾಣ್ಯ} (nāṇya) — a coin}

\label{lesson:KA-C185-nanya}



\begin{quote}
Before the new one: say the Kannada for tender coconut water, then the Kannada for a banknote.
\end{quote}

\subsection*{You'll want to know: \kn{ನಾಣ್ಯ}}
\textbf{\kn{ನಾಣ್ಯ}} — \emph{nāṇya} — \textquotedblleft{}a coin\textquotedblright{}.

\begin{grammarlens}[title={What you've built}]
One more word for this chapter.
\end{grammarlens}

\subsection*{Your turn}
\begin{itemize}
\raggedright
  \item \emph{Say it:} \emph{nāṇya}
  \item \emph{Say it:} \emph{nāṇya}, once more
  \item \emph{Say it:} say \emph{nāṇya}
  \item \emph{Recall:} say the Kannada for tender coconut water, then the Kannada for a banknote, then say \emph{nāṇya} again
\end{itemize}

\begin{tcolorbox}[breakable,colback=teal!4,colframe=teal!35!black,title={Before you move on}]
What does \kn{ನಾಣ್ಯ} mean? (\textquotedblleft{}A coin\textquotedblright{}.) Say it once more.
\end{tcolorbox}
\section[\texorpdfstring{\kn{ರಿಯಾಯಿತಿ} (riyāyiti)}{ರಿಯಾಯಿತಿ (riyāyiti)}]{\kn{ರಿಯಾಯಿತಿ} (riyāyiti) — a discount}

\label{lesson:KA-C185-riyayiti}



\begin{quote}
Before the new one: say the Kannada for a banknote, then the Kannada for a coin.
\end{quote}

\subsection*{You'll want to know: \kn{ರಿಯಾಯಿತಿ}}
\textbf{\kn{ರಿಯಾಯಿತಿ}} — \emph{riyāyiti} — \textquotedblleft{}a discount\textquotedblright{}.

\begin{grammarlens}[title={What you've built}]
One more word for this chapter.
\end{grammarlens}

\subsection*{Your turn}
\begin{itemize}
\raggedright
  \item \emph{Say it:} \emph{riyāyiti}
  \item \emph{Say it:} \emph{riyāyiti}, once more
  \item \emph{Say it:} say \emph{riyāyiti}
  \item \emph{Recall:} say the Kannada for a banknote, then the Kannada for a coin, then say \emph{riyāyiti} again
\end{itemize}

\begin{tcolorbox}[breakable,colback=teal!4,colframe=teal!35!black,title={Before you move on}]
What does \kn{ರಿಯಾಯಿತಿ} mean? (\textquotedblleft{}A discount\textquotedblright{}.) Say it once more.
\end{tcolorbox}
\section[\texorpdfstring{\kn{ಕಿಲೋ} (kilō)}{ಕಿಲೋ (kilō)}]{\kn{ಕಿಲೋ} (kilō) — a kilogram}

\label{lesson:KA-C185-kilo}



\begin{quote}
Before the new one: say the Kannada for a coin, then the Kannada for a discount.
\end{quote}

\subsection*{You'll want to know: \kn{ಕಿಲೋ}}
\textbf{\kn{ಕಿಲೋ}} — \emph{kilō} — \textquotedblleft{}a kilogram\textquotedblright{}.

\begin{grammarlens}[title={What you've built}]
One more word for this chapter.
\end{grammarlens}

\subsection*{Your turn}
\begin{itemize}
\raggedright
  \item \emph{Say it:} \emph{kilō}
  \item \emph{Say it:} \emph{kilō}, once more
  \item \emph{Say it:} say \emph{kilō}
  \item \emph{Recall:} say the Kannada for a coin, then the Kannada for a discount, then say \emph{kilō} again
\end{itemize}

\begin{tcolorbox}[breakable,colback=teal!4,colframe=teal!35!black,title={Before you move on}]
What does \kn{ಕಿಲೋ} mean? (\textquotedblleft{}A kilogram\textquotedblright{}.) Say it once more.
\end{tcolorbox}
\section[\texorpdfstring{\kn{ಪರ್ಸ್} (pars)}{ಪರ್ಸ್ (pars)}]{\kn{ಪರ್ಸ್} (pars) — a purse, a wallet}

\label{lesson:KA-C185-pars}



\begin{quote}
Before the new one: say the Kannada for a discount, then the Kannada for a kilogram.
\end{quote}

\subsection*{You'll want to know: \kn{ಪರ್ಸ್}}
\textbf{\kn{ಪರ್ಸ್}} — \emph{pars} — \textquotedblleft{}a purse, a wallet\textquotedblright{}.

\begin{grammarlens}[title={What you've built}]
One more word for this chapter.
\end{grammarlens}

\subsection*{Your turn}
\begin{itemize}
\raggedright
  \item \emph{Say it:} \emph{pars}
  \item \emph{Say it:} \emph{pars}, once more
  \item \emph{Say it:} say \emph{pars}
  \item \emph{Recall:} say the Kannada for a discount, then the Kannada for a kilogram, then say \emph{pars} again
\end{itemize}

\begin{tcolorbox}[breakable,colback=teal!4,colframe=teal!35!black,title={Before you move on}]
What does \kn{ಪರ್ಸ್} mean? (\textquotedblleft{}A purse, a wallet\textquotedblright{}.) Say it once more.
\end{tcolorbox}
